\dossier{4}{Iunianus et Victora, « amis des pauvres »}
\trad{Déposition de Iunianus, la veille des ides d’avril (12 avril), sous le consulat de Marcellinus et de Probinus. Il vécut quarante ans et mourut en paix ; ami des pauvres, il vécut avec sa \emph{virginia} quinze ans. À lui, bien méritant, sa \emph{virginia} Victora (\emph{Bictora}), bien méritante, a fait (ce monument) : amie des pauvres et travailleuse.}
\comm{La plus ancienne épitaphe du corpus qui emploie \emph{virginia} pour l’épouse épousée vierge (chapitre~\ref{ch:personnes}). Les deux époux reçoivent la même qualification d’« amis des pauvres » (\emph{amator}, \emph{amatrix pauperum}, avec la graphie \emph{paupe\{ro\}rum} que l’EDCS corrige et que l’EDB conserve) ; l’épouse est en outre \emph{operaria}, mot qui peut désigner une travailleuse manuelle ou, dans un sens plus particulier, une personne active dans les œuvres charitables, ce que le texte ne permet pas de préciser. Le bêtacisme est constant (\emph{bixit}, \emph{birginia}, \emph{Bictora}). L’EDB décrit un monogramme du Christ à la fin du texte, que l’EDCS ne signale pas. La date de 341 repose sur une formule consulaire lue en entier.}

\dossier{15}{Cassia, \emph{clarissima femina}}
\trad{[…] Cassia, (fille ou épouse) de Piso, \emph{clarissima femina}, vécut trente ans et quatre-vingt-treize jours ; née sous le consulat de Sabinus et Rufinus, le dix-septième jour avant les calendes de juillet (15 juin 316), elle rendit (l’âme) après le consulat d’Amantius et d’Albinus, le dix-septième (EDB : seizième) jour avant les calendes d’octobre (15 ou 16 septembre 346).}
\comm{Le génitif \emph{Pisonis} peut désigner le père ou le mari. Le texte est gravé sur un sarcophage (17 × 130 cm selon l’EDB, qui donne sans doute les dimensions du seul champ inscrit). Il contient deux dates consulaires, celle de la naissance (316) et celle de la mort (post-consulat de 345, soit 346) ; l’année de classement est celle de la mort. La durée de vie gravée, trente ans et quatre-vingt-treize jours, concorde avec les deux dates, selon le mode de décompte (chapitre~\ref{ch:calendrier}) : du 15 juin au 15 septembre, on compte quatre-vingt-treize jours en décompte inclusif, et autant jusqu’au 16 septembre en décompte exclusif. Les deux bases divergent sur le quantième du décès (XVII ou XVI), et l’EDB restitue le nom du mois ; la concordance arithmétique ne permet pas de choisir entre les deux lectures. Le verbe \emph{reddit}, employé absolument, appartient au vocabulaire de la mort comme restitution (de l’âme ou de la vie).}

\dossier{22}{Caianus, juriste}
\trad{(Vers) Heureuse fut la vie de cet homme, plus heureux encore son trépas : le nom de Caianus croît sans cesse au fil des siècles ; la renommée ne sait pas mourir, mais vit toujours et partout. Hôte romain, il vint dans la ville du prince ; il fut ici d’abord jurisconsulte et ami. — Il repose en paix ; déposé le quatrième jour avant les nones d’août (2 août), sous le consulat de Flavius Philippus et de Flavius Salia (348). Son père Sabbatius a fait (ce monument).}
\comm{Épitaphe métrique (CLE 659) d’un jurisconsulte. Les quatre hexamètres reprennent les lieux communs de la gloire posthume (\emph{nescit fama mori}) ; la partie en prose emploie le formulaire chrétien (\emph{quiescit in pace}, avec la graphie \emph{pacem}). Le vers \emph{advenit hospes Romanus princeps in urbem} est obscur : \emph{princeps} peut qualifier la ville (« la ville capitale ») ou désigner le prince ; la traduction proposée reste incertaine. L’EDB décrit un oiseau et un chrisme radié. Le père porte un nom dérivé du sabbat, \emph{Sabbatius} (chapitre~\ref{ch:noms}).}

\dossier{26}{Esuvia Ianuaria et Aradius Melissus}
\trad{D’une bonté et d’une sainteté admirables : à Esuvia Ianuaria, son épouse, qui vécut avec moi dix-huit ans, cinq mois et vingt-cinq jours, Aradius Melissus, son mari, a fait (ce monument), et pour lui-même ; déposée en paix le troisième jour avant les calendes d’avril (30 mars), sous le consulat de Philippus et de Salia (348).}
\comm{Le texte ne donne pas l’âge de la défunte, mais la durée de la vie commune (\emph{vixit mecum}), avec la même précision qu’un âge ; la table des personnes ne lui attribue donc pas d’âge (chapitre~\ref{ch:personnes}). Le mari prépare sa propre sépulture (\emph{et sibi fecit}). Le texte n’est pas apparié à l’EDB ; il est connu par l’ICVR II 5711 et par une publication de 2012.}

\dossier{36}{Clodius Adelfius, ancien préfet de la Ville}
\trad{Clodius Adelfius, \emph{vir clarissimus}, ancien préfet de la Ville, à son épouse incomparable et pour lui-même a fait (ce monument).}
\comm{Le texte est gravé en colonne étroite, une ou deux syllabes par ligne, disposition que les deux bases reproduisent. Il ne contient aucune date : l’année 351 de la base dépend de l’identification de Clodius Adelfius à Clodius Celsinus Adelphius, préfet de la Ville en 351, et marque un terme postérieur, non l’année du décès (chapitre~\ref{ch:fondement}). La bibliographie fait de ce préfet l’époux de la poétesse chrétienne Faltonia Betitia Proba ; l’identification et cette relation sont à contrôler dans la \emph{Prosopography of the Later Roman Empire}. Le texte, sans formule chrétienne, est classé dans le corpus en raison de sa publication parmi les inscriptions chrétiennes (ICVR I 19).}

\dossier{38}{Privata, \emph{in pace} et le monogramme}
\trad{Privata, qui vécut cinquante-trois ans, huit mois et deux jours, a été déposée le quatrième jour avant les nones de février (2 février) ; bien méritante, en paix ☧ ; sous le consulat de Decentius et de Paulus (352).}
\comm{Dossier de référence pour la question du monogramme (chapitre~\ref{ch:dieu_christ}). L’EDCS développe \emph{Chr(isti)} ; l’EDB transcrit \emph{((monogramma Christi))} ; l’estampage conservé à l’Ohio State University sous le numéro ILCV 3252A montre un christogramme placé après IN PACE. La pierre porte donc un signe, non le mot écrit en lettres. La traduction rend le signe par ☧ sans décider s’il se lit \emph{Christi} (« dans la paix du Christ ») ou s’il est juxtaposé à la formule. Le texte provient, selon l’EDB, du cimetière à ciel ouvert proche de la basilique Saint-Pierre ; il est conservé au Lapidario Cristiano des Musées du Vatican. La forme \emph{meses} pour \emph{menses} est signalée par les deux bases.}

\dossier{43}{Anastasius, poème de Constantia}
\trad{(Vers) Au nom [du Christ …]. Affligée, [Constantia écrit ces vers] pour Anastasius, son époux, qui [a échangé] la lumière [contre d’amères ténèbres] ; sa vie [s’arrêta] à quatre fois dix [et cinq ans] ; combien vite il fut ravi [à l’amour de son épouse chérie] ! Pleuré [quand Janus prenait] son douzième [lever], il fut enseveli [l’année que conduisait] le consul Arbitio (355).}
\comm{Les deux bases donnent la même restitution, empruntée à l’édition : chaque vers est conservé sur sa moitié gauche, et la seconde moitié est reconstituée selon le mètre et le sens. Le nom du Christ, dans l’invocation initiale, est entièrement restitué ; l’EDB décrit un chrisme avec alpha et oméga, qui rend l’invocation vraisemblable sans la rendre lisible. L’année repose sur le nom d’Arbitio, lu (\emph{Arbitio co[nsul]}) ; la date du jour (\emph{duodecimum [Ianus]}) est une restitution. L’âge, « quatre fois dix [et cinq] », est en partie restitué.}

\dossier{70}{Secunda, \emph{virgo fidelis}}
\trad{À Secunda, d’une bonté admirable, qui vécut d’une foi pure vingt ans ; chaste, elle a cessé (de vivre) en paix, vierge fidèle ; bien méritante, elle repose ; aux ides de juillet (15 juillet) ; colombe sans fiel ; sous (le consulat de) M(amertinus) et N(evitta) (362).}
\comm{L’un des deux seuls textes datés par les seules initiales des consuls (chapitre~\ref{ch:fondement}) ; l’identification de l’année 362 repose sur le couple M. et N., que seul le collège de Mamertinus et Nevitta satisfait dans la période. Le vocabulaire est celui de la virginité consacrée probable (\emph{pura fide}, \emph{pudica}, \emph{virgo fidelis}) ; l’image de la colombe sans fiel (\emph{palumbo sine felle}, au masculin) appartient au langage chrétien de l’innocence. L’EDB lit \emph{quiescet} et conserve \emph{\{id\}}, que l’EDCS considère également comme superflu.}

\dossier{72}{Flavia Balbilla : \emph{Dis Manibus} et \emph{in pace}}
\trad{Consacré aux dieux Mânes. Flavia Balbilla vécut vingt-huit ans, six mois et douze jours ; elle a reposé en paix sous le consulat de Mamertinus et de Nevitta, Maximus étant préfet de la Ville, le cinquième jour avant les calendes de février (28 janvier 362).}
\comm{Seul texte de l’ensemble C qui porte la dédicace aux Mânes, associée au verbe chrétien du repos (\emph{requievit in pace}). Les deux formules coexistent sans difficulté apparente pour le commanditaire ; la dédicace \emph{D. M. s.} est alors une formule d’ouverture traditionnelle plus qu’une profession religieuse. Le texte est aussi le seul du corpus à dater par le préfet de la Ville en plus des consuls (chapitre~\ref{ch:rangs}). L’EDB transmet la graphie \emph{Nebidda}, que l’EDCS normalise.}

\dossier{73}{Eleuteria : l’esprit, les saints et les hauteurs}
\trad{(Vers) [Voici] la demeure éternelle où tu reposes désormais en sécurité ; […] ton esprit, se retirant de la chair, [vit avec / est associé aux] saints, pour tes mérites et tes si grandes œuvres ; [toi qui …] tu reposes toujours en sécurité ; [tu as confié ton corps] à la terre et le gage reçu ; […] à la splendeur avec une claire lumière ; […] toi qui t’es toujours consacrée à Dieu ; […] mais pour tes actions tu es appelée vers les hauteurs ; […] Pacatianus doit beaucoup […]. Pacatianus a fait cet emplacement pour Eleuteria, sa mère, qui vécut environ soixante-quinze ans, deux mois et dix jours ; déposée le quatrième jour avant les calendes de juin (29 mai) ; elle dort [en paix], sous le consulat de notre seigneur Claudius Iulianus Auguste pour la quatrième fois et de Flavius Sallustius (363). Hesperius, qu’elle a nourri, l’a gravé.}
\comm{Le poème est conservé sur sa partie droite, et les deux bases restituent différemment le début des vers (chapitre~\ref{ch:typologie}). La traduction ne retient que ce qui est commun : l’esprit se sépare de la chair (lu), il est en relation avec les saints (\emph{sanctis}, lu ; le verbe est restitué), la défunte s’est consacrée à Dieu (\emph{deo dica⟨vi⟩sti}, lu dans l’EDB, \emph{dicasti} dans l’EDCS) et elle est appelée « vers les hauteurs » (lu). Le séjour \emph{in caelo} et la crainte de Dieu (\emph{[deu]m metuisti}) ne sont que des restitutions de l’EDCS ; l’EDB lit \emph{[- - - de]metuisti}. Le texte nomme trois personnes : la défunte, son fils Pacatianus, qui fait l’emplacement, et Hesperius, son nourricier, qui grave. L’EDB décrit des palmes, un soleil, un chrisme, une croix gammée, un alpha et un oméga, disposés à la fin des vers.}

\dossier{110}{Benenatus, ancien \emph{comes}, et Theophila}
\trad{(Grec) En paix a reposé Benenatus, ancien \emph{comes}, ayant vécu soixante-quinze ans ; il gît en ce lieu, sorti de la vie le quatrième jour avant les ides de mars (12 mars), sous le consulat de Lupicinus et de Iovinus (367) ; moi, Theophila, j’ai fait (ce monument) pour mon époux. — (Latin) À Benenatus, Theophila, son épouse, a fait (ce monument).}
\comm{Épitaphe bilingue, dont le grec porte l’essentiel (chapitre~\ref{ch:grec}). Le formulaire grec (\emph{ἐν εἰρήνῃ ἀνεπάη}) est l’équivalent du \emph{requiescit in pace} latin ; la date romaine est transposée en grec. L’abréviation \emph{ἀπ(ὸ) κο(μήτων)} désigne un ancien \emph{comes}, sans préciser la fonction. Les graphies \emph{κῖτε} et \emph{ἐξελτών} reflètent la prononciation contemporaine. Le nom de l’épouse, \emph{Theophila}, est théophore.}

\dossier{118}{Ianuaria : l’âge à l’heure près}
\trad{Sous le consulat de Lupicinus et de Iovinus (367), Ianuaria, qui vécut vingt-quatre ans, un jour et deux heures et demie, s’en est allée un mardi.}
\comm{Seul texte du corpus qui donne, sans restitution, une durée de vie en heures (chapitre~\ref{ch:personnes}). L’EDB ne transcrit pas \emph{s(emis)} et lit \emph{XXI III}, ce qui pourrait signaler une disposition particulière des chiffres sur la pierre ; elle décrit l’image de la défunte et des colombes. Le quantième n’est pas donné : le mardi ne peut être contrôlé (chapitre~\ref{ch:calendrier}). Les deux noms consulaires sont gravés au datif ou à l’ablatif avec \emph{in consulatu}, construction que l’EDCS signale par \emph{(!)}.}

\dossier{126}{Un fils né en 350, mort en 368}
\trad{Ses parents affligés ont fait ce que lui-même aurait dû faire ; malheureux parents, (contraints) par la douleur à ce devoir, ils ont fait aussi ce tombeau. Il naquit sous le consulat de Flavius Sergius et de Flavius Nigrinianus (350), le troisième jour avant les ides d’août (11 août), un samedi ; il a été déposé dans cette demeure éternelle sous le consulat de nos seigneurs Valentinien et Valens, Augustes pour la deuxième fois (368), le septième jour avant les ides de décembre (7 décembre), un samedi ; il vécut dix-neuf ans, trois mois et vingt-six jours ; en paix.}
\comm{Le topos de l’inversion des devoirs (les parents font pour l’enfant ce que l’enfant aurait dû faire pour eux) ouvre le texte. Les deux dates sont complètes. Le calcul confirme le samedi de la naissance ; le 7 décembre 368 est un dimanche, soit un écart d’un jour avec le samedi gravé. Les mois et les jours de l’âge concordent exactement avec l’intervalle entre les deux dates, mais l’intervalle est de dix-huit ans et non de dix-neuf (chapitre~\ref{ch:calendrier}). La durée gravée, \emph{XVIIII}, est lue identiquement par les deux bases ; l’écart d’une année reste inexpliqué. La tombe est une \emph{domus aeterna}.}

\dossier{133}{Felix, « appelé à la sainte foi »}
\trad{Felix, appelé à la sainte foi, s’en est allé en paix ; un si grand amour et une si grande affection sont gardés pour lui par ses amis à jamais ; de son vivant, il fut un réconfort, miséricordieux, connu de tous. Agrippina a fait (ce monument) pour son très doux mari, avec qui elle vécut sans blessure de l’âme trois ans et dix mois ; il fut dans le siècle, tant qu’il vécut, trente-deux ans ; déposé le treizième jour avant les calendes de septembre (20 août), sous le consulat de Valentinien, très noble enfant, et de Victor (369).}
\comm{L’expression \emph{sanctae fidei vocitus} peut évoquer une conversion ou un baptême ; le texte ne précise pas lequel (chapitre~\ref{ch:foi_eglise}). La durée de la vie commune (trois ans et dix mois) avait été enregistrée comme un âge dans la table des personnes ; elle est corrigée (chapitre~\ref{ch:personnes}). L’éloge de Felix (\emph{solacius}, \emph{misericors}) relève de la charité ; l’expression \emph{fuit in saeculum} oppose la vie terrestre à ce qui suit. L’EDB rattache le texte à la basilique Saint-Pierre.}

\dossier{134}{Une dalle familiale : un prêtre, un diacre et leur sœur}
\trad{(a) Endormissement de Dulcitius, prêtre, aux calendes de février (1\textsuperscript{er} février), et de Gregorius, diacre, le treizième jour avant les calendes de mai (19 avril), et de Catafronia (EDB : Califronia), leur sœur, le cinquième jour avant les ides de mars (11 mars). (b) Déposition de Faustina le quatrième jour avant les calendes de novembre (29 octobre), sous le consulat de Valentinien, très noble enfant, et de Victor (369). (c) Endormissement d’Agapè le quatrième jour avant les nones (EDB : aux nones) d’octobre, en paix.}
\comm{La grande dalle (105 × 190 × 15 cm selon l’EDB), retrouvée en morceaux remployés dans un mur de la basilique des Saints-Nérée-et-Achillée, porte trois textes distincts, que l’EDB distingue (a, b, c). Seule la partie b est datée : l’année de classement ne vaut pas pour les autres défunts (chapitre~\ref{ch:evenements}). Le texte atteste une famille cléricale (un prêtre, un diacre et leur sœur) et l’emploi de \emph{dormitio} pour désigner la mort. Les deux bases divergent sur le nom de la sœur et sur la date d’Agapè.}

\dossier{140}{Bonifatia, \emph{in pace Chr(isti)}}
\trad{À Bonifatia, bien méritante, dans la paix du Christ, qui vécut vingt-huit ans ; déposée en paix le douzième jour avant les calendes de septembre (21 août), sous le consulat de Valentinien et de Valens pour la troisième fois (370).}
\comm{La formule \emph{in pace Chr(isti)} suit \emph{bene merenti} et précède une seconde mention de la paix (\emph{deposita in pace}). Les deux bases transcrivent une abréviation, mais l’EDB décrit un chrisme parmi les données matérielles : il est probable que l’abréviation développée rende un monogramme employé comme compendium (chapitre~\ref{ch:dieu_christ}). L’EDB rattache le texte au cimetière de Pontien.}
